% Doplnky do kapitoly 1 bakalarskej prace
% Navrhnute ako bloky na vlozenie do existujucej struktury.

\section*{Doplnky do teoretickej casti}

\subsection*{1.2.6 Zhrnutie}
V prostredí platformy .NET existuje viacero prístupov k tvorbe webových aplikácií, pričom každý z nich prináša odlišný kompromis medzi komplexitou riešenia, interaktivitou používateľského rozhrania a mierou kontroly nad aplikačnou logikou. Z hľadiska navrhovaného systému sa ako najvhodnejší javí server-renderovaný prístup postavený na technológii ASP.NET Core MVC. Dôvodom je najmä prirodzená podpora formulárových a transakčných scenárov, integrované bezpečnostné mechanizmy, jednoduchšie nasadenie a vhodnosť pre aplikácie, ktoré potrebujú spoľahlivo spracovať autentifikáciu, autorizáciu a validáciu už na strane servera \cite{aspnetcore_overview}.

Pre navrhovaný elektronický obchod so zbraňami je dôležité, že server-renderovaný prístup umožňuje aplikovať obchodné a bezpečnostné pravidlá ešte pred odoslaním odpovede klientovi. Táto vlastnosť je významná najmä pri scenároch, v ktorých je potrebné rozlišovať medzi verejne dostupným sortimentom a regulovaným tovarom, kontrolovať oprávnenie používateľa alebo zaznamenávať auditné informácie o spracovaní objednávky. Zároveň ide o prístup, ktorý vhodne podporuje viacvrstvovú architektúru a oddelenie prezentačnej, aplikačnej a dátovej vrstvy.

Na základe uvedeného porovnania bol pre implementáciu systému zvolený framework ASP.NET Core MVC na platforme .NET 8. Tento výber zodpovedá charakteru navrhovanej aplikácie, ktorá kladie dôraz na bezpečný objednávkový proces, jednoznačne definované role používateľov, auditovateľnosť operácií a integráciu s databázovou a autentifikačnou vrstvou.

\subsection*{1.3.5 Zhrnutie}
Z pohľadu navrhovaného systému je rozhodujúce, že relačný model databázy umožňuje spoľahlivo reprezentovať vzťahy medzi používateľmi, produktmi, objednávkami, položkami objednávok a auditnými záznamami. Pre implementáciu prototypu bol preto ako primárny prístup zvolený Entity Framework Core, ktorý poskytuje dostatočnú úroveň abstrakcie, podporu migrácií a úzku integráciu s frameworkom ASP.NET Core \cite{efcore_overview}. Micro-ORM prístup je možné v budúcnosti využiť v špecifických čítacích alebo reportovacích scenároch, avšak pre jadro systému, ktoré pracuje s väzbami medzi entitami a so stavovými prechodmi objednávok, je vhodnejšie použiť plnohodnotný ORM nástroj.

Toto rozhodnutie sa priamo premieta do návrhu implementovaného systému, v ktorom databázová vrstva zabezpečuje perzistenciu produktového katalógu, používateľských účtov, objednávok, položiek objednávok, notifikácií a auditných záznamov. Dôležitou požiadavkou je aj podpora evolúcie databázovej schémy počas vývoja, ktorú EF Core rieši prostredníctvom migrácií.

\subsection*{1.4.2 Open-source a .NET e-commerce platformy}
V prostredí .NET patrí medzi najvýznamnejšie open-source e-commerce riešenia platforma nopCommerce. Ide o systém postavený na technológii ASP.NET Core, ktorý poskytuje rozsiahlu základnú funkcionalitu vrátane správy katalógu, objednávok, zákazníkov, SEO nastavení, platobných a dopravných modulov a podpory viacerých obchodov \cite{nopcommerce_features}. Výhodou tohto riešenia je technologická blízkosť k platforme .NET a možnosť modifikácie zdrojového kódu. Z pohľadu navrhovaného systému ide o dôležitý referenčný príklad, pretože potvrdzuje, že .NET platforma je plnohodnotne použiteľná pre implementáciu e-shop riešenia.

Na druhej strane, nopCommerce je všeobecná e-commerce platforma. Hoci podporuje rozšírenia a zásahy do kódu, implementácia doménovo špecifických pravidiel, ako je rozdelenie sortimentu na verejne dostupné položky a regulované zbrane, viacstupňové interné schvaľovanie objednávok alebo väzba na overovanie dokladov zákazníka, by si vyžadovala výraznejšie zásahy do návrhu systému. Z tohto pohľadu je nopCommerce vhodným referenčným riešením, ale nie optimálnym základom pre prototyp zameraný na predaj zbraní.

\subsection*{1.4.3 Komerčné a SaaS riešenia}
Významnú skupinu existujúcich riešení predstavujú SaaS a komerčné platformy, medzi ktoré patria napríklad Shopify alebo enterprise CMS systémy. Shopify poskytuje online obchod ako službu a umožňuje rýchle vytvorenie katalógu, vlastných menu, filtrov, responzívnej prezentácie produktov a urýchleného checkout procesu \cite{shopify_online_store}. Takýto model je vhodný pre štandardný online predaj, kde je prioritou rýchlosť nasadenia a minimalizácia technickej správy infraštruktúry.

Nevýhodou z pohľadu riešenej problematiky je závislosť od poskytovateľa a obmedzená možnosť zásahu do jadra obchodnej logiky. Pri regulovanom predaji, kde je potrebné implementovať vlastné validačné pravidlá, viacero interných rolí a auditovateľný workflow objednávky, môže byť tento model príliš reštriktívny. Komerčné CMS riešenia typu Kentico alebo Optimizely síce ponúkajú vyššiu mieru rozšíriteľnosti a profesionálnu podporu, avšak pre rozsah bakalárskej práce sú zbytočne robustné a licenčne náročné.

\subsection*{1.4.4 Open-source CMS ekosystém}
Osobitnú skupinu riešení predstavujú platformy vychádzajúce z CMS ekosystémov, napríklad WooCommerce nad WordPressom. WooCommerce umožňuje predávať fyzické aj digitálne produkty, podporuje rozšírenia, viacjazyčný obsah, rôzne meny a otvorené API rozhrania \cite{woocommerce_features}. Výhodou je najmä flexibilita, veľký ekosystém rozšírení a nízka vstupná bariéra pri budovaní menších obchodov.

Zároveň však ide o riešenie naviazané na WordPress, čo znamená odlišný technologický základ než platforma .NET. Pre bakalársku prácu orientovanú na návrh a implementáciu viacvrstvovej .NET aplikácie by použitie WooCommerce neprinášalo dostatočný priestor na demonštráciu architektonických rozhodnutí a implementácie vlastnej obchodnej logiky. WooCommerce je preto vhodné ako porovnávací referenčný bod, nie ako priamy technologický základ navrhovaného systému.

\subsection*{1.4.7 Zhrnutie}
Rešerš existujúcich riešení ukazuje, že univerzálne e-commerce platformy poskytujú rozsiahlu funkcionalitu pre bežný online predaj, avšak pri špecifickej doméne regulovaného predaja narážajú na limity v oblasti prispôsobenia procesov, bezpečnostných pravidiel a auditovateľnosti. nopCommerce predstavuje technologicky blízke open-source riešenie v prostredí .NET, Shopify demonštruje silu SaaS prístupu a WooCommerce flexibilitu open-source CMS ekosystému \cite{nopcommerce_features,shopify_online_store,woocommerce_features}.

Pre navrhovaný systém z toho vyplýva, že vlastná implementácia je opodstatnená. Umožňuje navrhnúť katalóg rozdelený na verejne dostupné položky a regulovaný tovar, implementovať overovanie zákazníckeho profilu, viacstupňové spracovanie objednávok internými rolami a viesť auditné záznamy o spracovaní objednávky. Práve tieto požiadavky predstavujú rozdiel medzi bežným e-shopom a špecializovaným systémom zameraným na predaj zbraní.

\subsection*{1.5 Metodická poznámka}
Táto kapitola vychádza z právneho stavu, z ktorého bol odvodený návrh prototypu a jeho validačné pravidlá. Cieľom kapitoly nie je podať úplný právny výklad, ale identifikovať tie legislatívne požiadavky, ktoré majú priamy vplyv na návrh dátového modelu, používateľských rolí a objednávkového procesu systému. Z tohto dôvodu je legislatíva interpretovaná predovšetkým z pohľadu požiadaviek na návrh a implementáciu aplikácie.

\subsection*{1.5.4 Zhrnutie legislatívnych požiadaviek na návrh aplikácie}
Z pohľadu implementácie systému znamenajú uvedené legislatívne obmedzenia potrebu rozdelenia katalógu na verejne dostupný sortiment a sortiment vyžadujúci dodatočné overenie. Aplikácia musí vedieť pracovať s údajmi o veku zákazníka, stave overenia profilu a s nahranými dokladmi, pričom samotná objednávka regulovaného tovaru nesmie byť dokončená bez splnenia definovaných podmienok. Rovnako dôležitá je evidencia objednávkového procesu, aby bolo možné spätne preukázať, kto a v akom stave objednávku schválil alebo spracoval.

Teoretické požiadavky legislatívnej časti sa tak v návrhu systému premietajú do validačných pravidiel, role-based autorizácie a do workflow objednávky, ktoré odlišuje verejne dostupné položky od zbraní viazaných na dodatočné obmedzenia. Tým sa legislatívna časť nestáva len opisom právneho rámca, ale priamo ovplyvňuje architektúru navrhovanej aplikácie.

\subsection*{1.6.5 Zhrnutie}
Z teoretických východísk bezpečnosti vyplývajú pre navrhovaný systém konkrétne implementačné požiadavky. Ide najmä o potrebu role-based autorizácie, server-side validácie vstupov, oddelenia verejného a regulovaného katalógu, evidencie auditných záznamov o spracovaní objednávok a obmedzenia prístupu k citlivým údajom používateľov a ich dokladom. V prípade navrhovaného systému to znamená, že bezpečnosť nie je riešená iba na úrovni prihlasovania používateľa, ale aj na úrovni obchodného procesu, kde jednotlivé kroky spracovania objednávky vykonávajú špecifické interné roly systému.

Bezpečnostné princípy sa zároveň prejavujú v návrhu administratívnych a interných častí aplikácie, kde je potrebné oddeliť oprávnenia administrátora, skladníka a zbrojárskej roly. Takto navrhnutý systém lepšie reflektuje reálne procesy regulovaného predaja a umožňuje transparentné spracovanie objednávok od vytvorenia až po finálne dokončenie.

\subsection*{BibTeX záznamy}
\begin{verbatim}
@misc{aspnetcore_overview,
  author = {{Microsoft}},
  title = {Overview of ASP.NET Core},
  year = {2025},
  url = {https://learn.microsoft.com/aspnet/core/introduction-to-aspnet-core?view=aspnetcore-8.0},
  note = {[cit. 2026-04-10]}
}

@misc{efcore_overview,
  author = {{Microsoft}},
  title = {Overview of Entity Framework Core},
  year = {2024},
  url = {https://learn.microsoft.com/en-us/ef/core/},
  note = {[cit. 2026-04-10]}
}

@misc{nopcommerce_features,
  author = {{nopCommerce}},
  title = {Online Store Software Features},
  year = {2026},
  url = {https://www.nopcommerce.com/en/features},
  note = {[cit. 2026-04-10]}
}

@misc{shopify_online_store,
  author = {{Shopify}},
  title = {Online Store},
  year = {2026},
  url = {https://help.shopify.com/en/manual/online-store},
  note = {[cit. 2026-04-10]}
}

@misc{woocommerce_features,
  author = {{WooCommerce}},
  title = {WooCommerce Features},
  year = {2026},
  url = {https://woocommerce.com/woocommerce-features/},
  note = {[cit. 2026-04-10]}
}

@misc{gdpr_official,
  author = {{European Union}},
  title = {Regulation (EU) 2016/679 (GDPR)},
  year = {2016},
  url = {https://eur-lex.europa.eu/eli/2016/679/oj},
  note = {[cit. 2026-04-10]}
}
\end{verbatim}
